\section{MicroPython pre vstavané systémy}
\label{sec:micropython}

MicroPython je minimalistická implementácia jazyka Python určená pre obmedzené hardvérové platformy. Je navrhnutý tak, aby poskytoval interaktívne prostredie (REPL), priamy prístup k perifériám a rýchle prototypovanie pri zachovaní čitateľnej a expresívnej syntaxe.

\subsection*{Architektúra run‑time a pamäťový model}
MicroPython implementuje heap‑based správu pamäte s garbage collection (GC). Pre návrh stabilných aplikácií je potrebné pochopiť obmedzenia dynamickej alokácie a strategicky minimalizovať vytváranie objektov v kritických slučkách. Kľúčové sú tiež mechanizmy pre "frozen" moduly (vkladanie Python kódu do firmvéru) a native moduly pre urýchlenie kritických sekcií.

\subsection*{Výkonnostné optimalizácie}
Pre zvýšenie výkonu a deterministiky v MicroPythone sú k dispozícii viaceré techniky:
\begin{itemize}
	\item Použitie dekorátora \texttt{micropython.native} alebo \texttt{micropython.viper} pre kompiláciu kritických funkcií do natívneho kódu.
	\item Vkladanie natívnych modulov v C pre časovo náročné operácie alebo pri práci s veľkými dátovými štruktúrami.
	\item Použitie DMA, hardvérových periférií a RMT namiesto softvérového bit‑bangingu.
\end{itemize}

\subsection*{Konkurentné modely a I/O}
MicroPython podporuje viaceré modely spracovania súbežných udalostí: prerušeniami riadené handlery (ISR), asynchrónne programovanie cez \texttt{uasyncio} a jednoduché preemptívne prístupy kombinované s frontami. ISR musia byť krátke a mali by len nastavovať príznaky alebo zapisovať do fronty, kde hlavný cyklus spracuje zložitejšie akcie.

\subsection*{Vývoj, nasadenie a nástroje}
Pre rýchle nahrávanie a ladenie sa často používa Thonny; pre produkčné buildy a integráciu sa odporúča automatisovať vytváranie firmware so zamrazenými modulmi a testovacím balíčkom. Dôležitý je aj monitoring pamäti počas vývoja a profilovanie časových kritických úsekov.

\subsection*{Silné a slabé stránky}
MicroPython ponúka výrazné zjednodušenie vývoja, interaktívne ladenie a rýchle nasadenie. Na druhej strane, oproti nativnému C/C++ riešeniu sú obmedzenia v rýchlosti, správe pamäte a dostupnosti knižníc. Preto pre produkčné systémy často kombinujeme MicroPython pre riadenie vyššej úrovne a natívne moduly pre výpočtovo náročné časti.

\subsection*{Praktické doporučenia}
\begin{itemize}
	\item Minimalizovať dynamickú alokáciu v hlavných slučkách a v kritických handleroch.
	\item Používať "frozen modules" pre stabilitu a zníženie nárokov na RAM.
	\item Využiť \texttt{uasyncio} pre čitateľné asynchrónne I/O a riadenie viacerých úloh.
	\item Keď je potrebná vysoká výkonnosť, extrahovať kritické algoritmy do C rozšírení.
\end{itemize}

\subsection*{Nasadenie}
MicroPython je vhodný pre prototypovanie, výučbu a menšie produkčné nasadenia, kde rýchlosť vývoja a flexibilita prevažujú nad maximálnou výkonnosťou. Pri väčších projektov sa odporúča hybridný prístup: MicroPython pre orchestráciu a natívne moduly pre časovo citlivé časti.

Pre ďalšie informácie a referencie pozri MicroPython dokumentáciu a príručky k platforme ESP32 \cite{micropython_docs}.
